\documentclass{article}
\usepackage[utf8]{inputenc}
\usepackage[portuguese]{babel}
\usepackage{amsmath}
\usepackage{amsfonts}
\usepackage{amssymb}
\usepackage{amsthm}

% Definições de ambientes matemáticos
%\newtheorem{proposition}{Proposição}
%\newtheorem{lemma}{Lema}
\newtheorem{corollary}{Corolário}

\usepackage{amsmath, amsthm, amssymb, booktabs, array}
\newtheorem{proposition}{Proposição}[section]
\newtheorem{lemma}[proposition]{Lema}
\newtheorem{remark}[proposition]{Observação}
\newtheorem{definition}[proposition]{Definição}

\title{Saturação Geométrica e Inevitabilidade Estrutural: Uma Abordagem de Malhas $3 \times N$ para a Conjectura de Goldbach}
\author{Tiago Bandeira \\ Demonstração de Estuturas em Grades $3 \times N$}
\date{Maio de 2026}

\begin{document}
	
	\maketitle
	
	\section{Introdução}
	Consideramos uma grade $G$ de dimensões $3 \times N$ preenchida por uma sequência de números ímpares $S = \{s_1, s_2, \dots, s_{3N}\}$. Definimos um mapeamento $f: S \to G$ tal que a distribuição de restos modulares não seja correlacionada com a geometria das colunas, garantindo a equiprobabilidade de ocupação para números primos em qualquer célula da grade.
	
	\section{Fundamentação Teórica}
	
	\begin{lemma}[Densidade de Primos em Ímpares]
		Seja $p(k)$ a probabilidade de um número ímpar $s_k$ ser primo. Pelo Teorema do Número Primo, a densidade de primos no conjunto dos ímpares é o dobro da densidade nos naturais para valores suficientemente grandes.
	\end{lemma}
	
	\begin{proof}
		Sabemos que $\pi(x) \approx \frac{x}{\ln x}$. Como os números ímpares representam metade do conjunto dos naturais, e todos os primos (exceto 2) são ímpares, a probabilidade local $p(k)$ para um termo $s_k$ é:
		\begin{equation}
			p(k) = \frac{P(\text{primo} \cap \text{ímpar})}{P(\text{ímpar})} = \frac{1/\ln(s_k)}{1/2} = \frac{2}{\ln(s_k)}
		\end{equation}
	\end{proof}
	
	\section{Demonstração da Proposição Principal}
	
	\begin{proposition}
		Em uma grade $3 \times N$ com mapeamento não-obstrutivo, a existência de ao menos uma diagonal ou coluna central composta exclusivamente por primos é um evento quase certo para $N \to \infty$.
	\end{proposition}
	
	\begin{proof}
		Seja $E_i$ o evento de sucesso na $i$-ésima janela $3 \times 3$ da grade. Para que uma diagonal ou a coluna central seja formada, três células específicas devem ser preenchidas por primos. Sob a hipótese de independência local de Cramér, a probabilidade de uma fileira é $p^3$. Usando o Princípio da Inclusão-Exclusão para as duas diagonais ($D_1, D_2$) e a vertical central ($C_v$):
		\begin{equation}
			P(E_i) = 3p(k)^3 - 3p(k)^5 + p(k)^7
		\end{equation}
		Para provar que o evento ocorre em algum ponto de $N$, analisamos a probabilidade complementar de falha em todas as janelas:
		\begin{equation}
			\mathcal{P}(\text{Falha Total}) = \prod_{i=1}^{N-2} (1 - P(E_i))
		\end{equation}
		Conforme $N \to \infty$, se a série $\sum P(E_i)$ divergir, o produto acima converge para zero.
	\end{proof}
	
	\begin{lemma}[Divergência da Série de Probabilidades]
		A soma das probabilidades de fileiras em janelas independentes diverge, garantindo a recorrência do evento.
	\end{lemma}
	
	\begin{proof}
		A probabilidade local é dominada pelo termo $p(k)^3$. Portanto, analisamos a série:
		\begin{equation}
			\sum_{k=1}^{\infty} \left( \frac{2}{\ln(s_k)} \right)^3 = 8 \sum_{k=1}^{\infty} \frac{1}{\ln(2k-1)^3}
		\end{equation}
		Pelo teste da integral, comparamos com $\int \frac{1}{(\ln x)^3} dx$. Como $\frac{1}{(\ln x)^3} > \frac{1}{x}$ para $x$ suficientemente grande, e $\int \frac{1}{x} dx$ diverge, a série de probabilidades também diverge. Pelo Segundo Lema de Borel-Cantelli, o evento $E$ ocorre infinitas vezes.
	\end{proof}

	
	\begin{corollary}
		A probabilidade de existência $\mathcal{P}_N$ tende a 1 conforme o comprimento $N$ da grade aumenta, independentemente da raridade crescente dos números primos.
	\end{corollary}
	
	\begin{proof}[Prova do Corolário]
		Decorre diretamente da aplicação do Segundo Lema de Borel-Cantelli sobre a divergência demonstrada no Lema 2.
	\end{proof}
	
	\begin{proposition}[Estimativa de Unicidade Estrutural]
		Seja $P(E_i)$ a probabilidade de ocorrência de uma configuração prima em uma janela $3 \times 3$ na coluna $i$. A expectativa do número de estruturas únicas $\mathcal{K}_u$ em uma grade $3 \times N$ é simplificada por:
		\begin{equation}
			\mathcal{K}_u \approx \frac{1}{3} \sum_{i=1}^{N-2} P(E_i)
		\end{equation}
		onde o fator $1/3$ compensa a redundância geométrica inerente ao deslocamento unitário da janela.
	\end{proposition}
	
	\begin{proof}
		Considerando que uma mesma diagonal ou coluna central de primos é detectada por três janelas sucessivas ($i, i+1, i+2$), a soma bruta de sucessos superestima a contagem real por um fator de 3. 
		Ao segmentarmos a grade em blocos de escala $a \gg 3$, a média local converge para a densidade teórica, e o erro de borda (estruturas que cruzam os blocos) torna-se desprezível, validando a proporcionalidade direta entre a soma das probabilidades locais e a existência de estruturas físicas inéditas.
	\end{proof}
	
	\vspace{0.5cm}
	\noindent \textbf{Observação (Definição da Probabilidade Local):} \\
	Vale ressaltar que a probabilidade $P(E_i)$ de uma configuração estrutural (como uma diagonal primária ou coluna central) se formar na janela $i$ é derivada do Teorema do Número Primo. Como a malha é construída exclusivamente por números ímpares, a densidade local de primos em torno de um valor $n_i$ é o dobro da densidade natural, ou seja, $P(p) \approx \frac{2}{\ln(n_i)}$. 
	
	Como a formação de uma estrutura geométrica de extensão 3 exige a primalidade simultânea de 3 células distintas, a probabilidade conjunta assumindo independência local é:
	\begin{equation}
		P(E_i) \approx \left( \frac{2}{\ln(n_i)} \right)^3 = \frac{8}{(\ln n_i)^3}
	\end{equation}
	onde $n_i$ representa o valor numérico médio estimado para a subgrade na coluna $i$. Ao substituir este termo na Equação (1), obtém-se o modelo computacional direto para a expectativa de unicidade.
	
	
	\begin{corollary}[Implicações de Ordem Inferior e a Conjectura de Goldbach]
		A convergência assintótica e a garantia de saturação estrutural de trios primos na malha $3 \times N$ implicam, por redução de complexidade probabilística, a inevitabilidade de partições de grau inferior. Este comportamento fornece forte sustentação analítica para proposições aditivas, notadamente a Conjectura de Goldbach.
	\end{corollary}
	
	\begin{proof}
		O modelo de unicidade para a malha geométrica exige a primalidade simultânea de três elementos dispostos em configurações rígidas (diagonais ou colunas centrais). A expectativa do número de estruturas, $\mathcal{K}_u$, possui uma taxa de crescimento da ordem de:
		\begin{equation}
			\mathcal{K}_u \propto \mathcal{O}\left(\frac{N}{(\ln N)^3}\right)
		\end{equation}
		
		Em contrapartida, a Conjectura de Goldbach postula a existência de pares de primos $(p_1, p_2)$ cuja soma resulte em $2N$. Pelo modelo heurístico de Hardy-Littlewood, o número esperado de tais partições, denotado por $\mathcal{G}(N)$, escala com a seguinte ordem:
		\begin{equation}
			\mathcal{G}(N) \propto \mathcal{O}\left(\frac{N}{(\ln N)^2}\right)
		\end{equation}
		
		Para valores suficientemente grandes de $N$, é evidente que a densidade de probabilidade obedece à relação $\frac{1}{(\ln N)^2} \gg \frac{1}{(\ln N)^3}$. A condição para a formação de pares aditivos é, portanto, consideravelmente menos restritiva do que a necessária para a formação geométrica de trios na malha. 
		
		Uma vez demonstrado que a distribuição normalizada dos primos assegura a divergência e o "transbordamento" de sucessos para o sistema restrito de grau 3 (trios) a partir de um limite de estabilidade (e.g., $N \ge 11.000$), conclui-se que o sistema engloba cobertura assintoticamente completa para interações de grau 2 (pares). Ao validar a estrutura tridimensional local, a malha atua concomitantemente como um gerador exaustivo de pares menores, atestando a robustez probabilística e a inevitabilidade das partições requeridas por Goldbach.
	\end{proof}
	
	
	
	
	
	% ============================================================
	%  Seção: Análise Assintótica de K_u e a Convergência de Malha
	% ============================================================
	
	\section{Análise Assintótica de $\mathcal{K}_u$ e o Limiar de Estabilidade}
	
	\subsection*{Notação e Contexto}
	
	Seja $N$ o tamanho do problema e $f(N)$ a contagem da sequência de ímpares $\{1, 3, 5, \ldots, N\}$, dada por $f(N) \sim N/2$. Consideramos a grandeza $\mathcal{K}_u$ como a densidade acumulada de eventos estruturais (sucessos) sem fatores de atenuação externa, visando observar o comportamento intrínseco da malha.
	
	% ------------------------------------------------------------------
	\begin{definition}[Densidade Acumulada $\mathcal{K}_u$]
		\label{def:ku_direta}
		A expectativa do número de eventos estruturais únicos em uma grade $3 \times N$, baseada na probabilidade local de trios primos, é definida por:
		\begin{equation}
			\mathcal{K}_u = \sum_{i=1}^{N-2} \frac{8}{(\ln n_i)^3}
			\label{eq:ku_direta}
		\end{equation}
		onde o numerador 8 decorre da densidade ajustada para números ímpares $(2^3)$.
	\end{definition}
	
	% ------------------------------------------------------------------
	\begin{lemma}[Crescimento Assintótico]
		\label{lem:ku_fast}
		A grandeza $\mathcal{K}_u$ satisfaz a seguinte equivalência quando $N \to \infty$:
		\begin{equation}
			\mathcal{K}_u \;\sim\; \frac{8N}{(\ln N)^3}.
		\end{equation}
	\end{lemma}
	
	% ------------------------------------------------------------------
	\begin{proposition}[Ponto de Virada e Saturação Estrutural]
		\label{prop:virada}
		Seja $\Phi(\mathcal{K}_u) = \frac{\mathcal{K}_u^2 + \mathcal{K}_u}{2}$ a função que descreve o acúmulo de interações estruturais. Para $N \geq 11.000$, a densidade de eventos supera o crescimento linear da malha $f(N) = N/2$. 
		Sob esta condição, tem-se:
		\begin{equation}
			\Phi(\mathcal{K}_u) \;\sim\; \frac{32\,N^2}{(\ln N)^6}
		\end{equation}
		com ponto de cruzamento empírico $N^* \approx 1{,}1 \times 10^4$.
	\end{proposition}
	
	\begin{proof}
		Pelo Lema~\ref{lem:ku_fast}, o termo dominante de $\Phi(\mathcal{K}_u)$ é $\mathcal{K}_u^2/2$. Comparando com $f(N)$:
		\[
		\frac{\Phi(\mathcal{K}_u)}{f(N)} \;\sim\; \frac{64\,N}{(\ln N)^6}
		\]
		Ao resolver a equação transcendental $64N = (\ln N)^6$ para o ponto de equilíbrio, encontramos:
		\[
		N^* \approx 11.000
		\]
		Neste limiar, a taxa de criação de novas estruturas únicas ultrapassa a taxa de crescimento dos elementos da malha, caracterizando o estado de saturação estrutural observado experimentalmente.
	\end{proof}
	
	% ------------------------------------------------------------------
	\begin{remark}[Modelo Direto vs. Estrito e a Dinâmica de Reciclagem]
		É fundamental distinguir o \textit{Modelo de Densidade Direta} do \textit{Modelo de Unicidade Estrita}. No modelo direto, o ponto de virada ocorre em $N^* \approx 1{,}1 \times 10^4$, alinhando-se à realidade computacional da malha. Este fenômeno é impulsionado pela "reciclagem" de primos pequenos e pela sobreposição de janelas, que aceleram a convergência e tornam a estrutura autossustentável precocemente.
		
		Caso seja aplicado o fator de correção de $1/3$ (Modelo Estrito) para eliminar redundâncias espaciais, o ponto de cruzamento desloca-se em duas ordens de magnitude ($N^* \approx 1{,}1 \times 10^6$). Ressalta-se, contudo, que o \textbf{comportamento assintótico} é invariante: em ambos os casos, $\Phi(\mathcal{K}_u)$ diverge como $\omega(N)$. A escolha pelo modelo direto em análises de partição (como Goldbach) justifica-se pela natureza cumulativa da densidade de primos, enquanto o modelo rígido oferece um limite superior de segurança teórica.
	\end{remark}
	
	
\end{document}